\chapter{总结与展望}


以上就是笔者在使用Latex总结的一些基本经验，可能也有一些不恰当的地方，欢迎评论和指正。Latex具体的一些操作实在有太多太多，大家也可以参考网上其他教程进行相关的高级操作，这里就不再赘述啦。

（最后一章为“总结与展望”。包括对整个论文主要成果的总结，本研究的创造性成果或创新点理论，对应用前景和社会、经济价值等的预测和评价，并指出今后进行进一步研究工作的展望与设想。写作时删除此段说明。）

这里引用两个参考文献\cite{2001Applying}\cite{2004PSO_ZhangLibiao}。

还可以这样引用\cite{2001Applying,2021A}。

或者说这样引用\cite{2001Applying,2004PSO_ZhangLibiao,2021A}

引用作者是这样引用\citet{2001Applying},\citet{2004PSO_ZhangLibiao}

或者是这样引用\citep{2001Applying}

本模板的参考文献采用国标 GB/T 7714-2015 的数字顺序样式（gbt7714-numeric），按论文中引用的先后顺序用阿拉伯数字连续编号。
如需改用作者-年份等其它排序方式，修改 \texttt{chapters/ref.tex} 中的参考文献样式设置即可。

在提交的论文中，有些同学出现了以下一些参考文献的错误，注意一些常见问题：
\begin{enumerate}
    \item 参考文献的写法不够统一，例如会议信息有的写了 Proceedings、有的没写。建议按官方提供的 bib 格式补全信息后再生成 PDF，会议名称等字段务必写全；
    \item 关于 arXiv 上的论文，常出现缺少期卷、页码的情况，注意按规范格式补全。arXiv 文件一般可归入 R、A、Z/OL 类别，有需要可自行酌情修改；
    \item 参考文献可以去网站DBLP\footnote{\url{https://dblp.org}}上进行检索，可以引入较为完整的bib格式，但是需要注意字段的修改。
\end{enumerate}
